\chapter{Le moteur de stratégies}\label{ch:ll:the-strategy-engine}

Exécutez deux fois la même stratégie de cotation sur la même journée enregistrée et comparez ce qu'elle a envoyé. Si un composant horodate ses ordres
par l'horloge murale, les deux exécutions ne concordent jamais~: cent relances donnent cent sorties différentes. Si un autre conserve ses cotations
dans un ensemble de hachage dont la graine est tirée au démarrage et les parcourt dans l'ordre de l'ensemble, 93 relances sur cent diffèrent. Si deux
threads alimentent sa file, 49 diffèrent. Supprimez les trois sources et les cent exécutions concordent jusqu'au dernier bit~; dès lors, tout bogue observé en
production peut être rejoué sur un poste de travail. Ce chapitre construit le composant qui fait tourner les stratégies de la société~: un thread unique
qui traite les événements de marché, les comptes rendus d'ordres, les temporisateurs et les changements de paramètres un à un, chacun jusqu'à complétion, dans un ordre
fixé par des règles et non par le système d'exploitation, avec un temps injecté plutôt que lu, et sans rien sur son chemin critique (hot path) qui
puisse attendre, allouer ou varier.

\section{La boucle d'événements}

\begin{definition}[Boucle d'événements, traitement jusqu'à complétion]\label{def:ll:the-strategy-engine:loop}
Une \emph{boucle d'événements}\index{boucle d'événements} est un thread qui prend de façon répétée l'événement suivant parmi ses entrées et le confie au code
qui le traite. Le \emph{traitement jusqu'à complétion}\index{traitement jusqu'à complétion} (run-to-completion) traite chaque événement entièrement, y compris
tout ce qu'il provoque, avant de prendre l'événement suivant~: aucun gestionnaire n'est interrompu par un autre, de sorte qu'un gestionnaire ne voit jamais un état
qu'un autre aurait modifié à moitié.
\end{definition}

Les entrées du moteur sont les anneaux du chapitre~12~: les événements de marché venant du \omterm{def:ll:the-feed-handler:handler}{feed handler} (par l'intermédiaire du \omterm{def:ll:the-order-book-builder:builder}{constructeur de carnet}
du chapitre~19), les comptes rendus de la passerelle d'ordres (chapitre~21) et les messages de contrôle~; ses propres temporisateurs proviennent d'une \omterm{def:ll:the-strategy-engine:wheel}{roue de temporisateurs}. Ses
sorties vont à l'anneau de la passerelle et au journal. Un seul thread possède tout l'état de la stratégie, si bien qu'aucun verrou n'est nécessaire, et
l'ordre de traitement des événements est fixé par une seule règle~: avant de traiter un événement de temps $t$, tous les temporisateurs échus à
$t$ ou avant se déclenchent, dans l'ordre des dates puis des identifiants~; les changements de paramètres prenant effet à $t$ ou avant sont appliqués~; puis
l'événement. La \Cref{fig:ll:the-strategy-engine:loop} en donne le schéma, et le \cref{lst:ll:the-strategy-engine:run} présente la boucle entière.

\begin{omfigure}
\begin{tikzpicture}[font=\footnotesize,box/.style={draw,rounded corners=2pt,minimum height=0.7cm,align=center,inner sep=3pt},
  >=Stealth]
\node[box,fill=omDef!15] (mkt) at (0,1.2) {événements de marché\\ (constructeur de carnet)};
\node[box,fill=omDef!15] (rep) at (0,0) {comptes rendus d'ordres\\ (passerelle)};
\node[box,fill=omDef!15] (ctl) at (0,-1.2) {instantanés\\ de paramètres};
\node[box,fill=omMeth!15] (wheel) at (3.6,-2.5) {roue de temporisateurs};
\node[box,fill=omThm!12,minimum width=4.2cm,minimum height=2.4cm] (loop) at (4.6,0) {};
\node[align=center] at (4.6,0.7) {\textbf{un seul thread}};
\node[align=left,font=\scriptsize,anchor=north] at (4.6,0.35) {1. temporisateurs échus à $t$\\ 2. paramètres effectifs à $t$\\ 3. l'événement en $t$};
\node[box,fill=omDef!15] (gw) at (9.0,0.6) {ordres\\ (anneau de la passerelle)};
\node[box,fill=gray!10] (log) at (9.0,-0.6) {journal};
\draw[->] (mkt) -- (loop);
\draw[->] (rep) -- (loop);
\draw[->] (ctl) -- (loop);
\draw[->] (wheel.north) -- (wheel.north |- loop.south);
\draw[->,dashed] ([xshift=0.9cm]loop.south) |- node[pos=0.25,right,font=\scriptsize]{planification} (wheel.east);
\draw[->] (loop) -- (gw);
\draw[->] (loop) -- (log);
\node[font=\scriptsize,align=left,anchor=west] at (7.4,-2.0) {horloge~: injectée\\ (simulée ou réelle)};
\end{tikzpicture}

\omcaption{Le moteur de stratégies. Toutes les entrées atteignent un seul thread, qui traite chaque événement jusqu'à complétion dans un ordre fixé par
une règle~: les temporisateurs échus à l'heure de l'événement, puis les paramètres en vigueur à ce moment, puis l'événement. Le temps est celui de l'horloge
donnée au moteur~: celui d'une journée enregistrée en rejeu, celui de la machine en temps réel.}\label{fig:ll:the-strategy-engine:loop}
\end{omfigure}

\omcode{code/firm/stratengine/cpp/firm_stratengine.hpp}{128}{136}{La boucle d'événements~: temporisateurs d'abord, puis paramètres, puis
l'événement, chacun traité jusqu'à complétion.}\label{lst:ll:the-strategy-engine:run}

La conception est ancienne et volontairement sobre. La plateforme de courtage dont Martin Fowler a décrit l'architecture en 2011 faisait tourner sa
logique métier sur un seul thread, en mémoire, avec un état «~entièrement dérivable du traitement des événements d'entrée~» conservé dans un
stockage durable~; et parce que «~la logique métier est déterministe~», ses sorties n'avaient même pas besoin d'être stockées. Un moteur de stratégies
en retire les deux mêmes avantages~: la vitesse, car un thread unique dont les données se trouvent dans ses propres caches ne se coordonne avec
personne, et le rejeu, car les mêmes entrées dans le même ordre produisent les mêmes sorties.

\section{Déterminisme monothread}

Un programme est reproductible bit à bit (la reproductibilité bit à bit du livre~4, chapitre~25) seulement dans son ensemble, et chacune
de ses parties peut le rompre. La \Cref{fig:ll:the-strategy-engine:reruns} chiffre les dégâts sur le petit moteur de cotation de la construction, exécuté cent fois comme un nouveau
processus sur les mêmes 6\,000 événements enregistrés, en n'activant qu'une source de non-déterminisme à la fois.

\begin{itemize}
\item \emph{Une lecture de l'horloge murale}, ici pour horodater chaque action et dater ses cotations~: chaque exécution lit des heures différentes, si bien que
cent exécutions donnent cent sorties. Toute décision qui compare le temps réel à quelque chose (l'âge d'une cotation, un délai d'expiration) fait
remonter la différence jusqu'aux ordres eux-mêmes.
\item \emph{Parcours d'un conteneur de hachage avec une graine propre au processus}~: le moteur conserve ses dix cotations dans un ensemble de hachage dont la fonction de hachage
est initialisée au démarrage et les recote dans l'ordre de l'ensemble, si bien que les identifiants et l'ordre des actions suivent la graine.
La table de hachage standard de Rust fait exactement cela par construction (son algorithme «~est initialisé aléatoirement~»)~; 93 relances sur cent
diffèrent.
\item \emph{Deux threads alimentant une même file}, chacun en portant la moitié des événements~: l'ordre dans lequel le consommateur les voit dépend
de l'ordonnancement, le carnet passe par des états qui n'ont jamais existé, et 49 relances sur cent diffèrent.
\end{itemize}

Désactivées, les trois sources donnent une seule sortie sur cent relances~; il en va de même pour le moteur de la société. Les règles qui le garantissent sont peu nombreuses. Le temps
provient des événements, par une \omterm{def:ll:the-strategy-engine:wheel}{horloge injectée} (section suivante). Les conteneurs parcourus ont un ordre défini~: tableaux,
tables ordonnées, ou tables de hachage dont l'ordre ne décide de rien. Un seul thread possède l'état~; les autres threads ne font que transmettre des données
par les anneaux, et la fusion de plusieurs entrées suit une règle (par temps, puis par source, puis par numéro de séquence), et non l'ordre
d'arrivée. Les sommes en virgule flottante se font dans un ordre fixé (chapitre~8), et la stratégie d'exemple reste en entiers, ce qui explique aussi que
les moteurs C++ et Rust de la construction produisent le même condensat de sortie.

\begin{omfigure}
\begin{tikzpicture}
\begin{axis}[width=11cm,height=5.4cm,ybar,bar width=16pt,ymin=0,ymax=115,
  symbolic x coords={firm,none,wall_clock,hash_order,two_threads,all},xtick=data,
  xticklabels={{moteur de\\ la société},{sans\\ défaut},{horloge\\ murale},{ordre de\\ hachage\\ à graine},{deux\\ threads},{les\\ trois}},
  xticklabel style={align=center,font=\footnotesize},enlarge x limits=0.1,ylabel={sorties distinctes sur 100 relances},
  tick label style={font=\small},label style={font=\small},grid=major,grid style={gray!25},
  nodes near coords,every node near coord/.append style={font=\scriptsize},table/col sep=comma]
\addplot[fill=omThm!45,draw=omThm] table[x=key,y=distinct]{figdata/low-latency/20-the-strategy-engine/measured_nondet.csv};
\end{axis}
\end{tikzpicture}

\omcaption{Nombre de condensats de sortie distincts sur cent relances (chacune un nouveau processus) des mêmes 6\,000 événements enregistrés, pour le moteur de la société
et pour le petit moteur de \texttt{ll\_nondet.hpp} avec chaque source de non-déterminisme activée. Mesuré sur un ordinateur portable
(Intel Core Ultra 7 155H) sous WSL2. Données~: \texttt{bench\_engine.py}.}\label{fig:ll:the-strategy-engine:reruns}
\end{omfigure}

\section{Temporisateurs et horloge injectée}

\begin{definition}[Roue de temporisateurs, horloge injectée]\label{def:ll:the-strategy-engine:wheel}
Une \emph{roue de temporisateurs}\index{roue de temporisateurs} est un tableau circulaire de cases, chacune couvrant un intervalle de temps (sa granularité)~: un
temporisateur est rangé dans la case de son instant d'échéance modulo l'étendue de la roue, et faire avancer la roue visite les cases du temps
écoulé et déclenche les temporisateurs qui y sont échus~; programmer un temporisateur coûte un temps constant quel que soit leur nombre. Une \emph{horloge
injectée}\index{horloge injectée} est une horloge passée au code qui a besoin du temps, au lieu d'être lue par lui~: le même code tourne
sur l'horloge simulée d'une journée enregistrée (livre~7, chapitre~17) ou sur l'horloge réelle de la machine.
\end{definition}

Une stratégie est pleine de temporisateurs~: un ordre qui doit être acquitté dans un délai, une cotation qui doit être retirée si les données de marché
s'arrêtent, un rafraîchissement toutes les quelques secondes, un \omterm{def:ll:protocols-i-fix:session}{heartbeat}. Une file de priorité suffirait, à un coût logarithmique par temporisateur~; la \omterm{def:ll:the-strategy-engine:wheel}{roue de temporisateurs} à hachage
de Varghese et Lauck (1987) le fait en temps constant, ce qui compte lorsque chaque ordre programme et annule des temporisateurs.
La roue du build compte 1\,024 cases de \qty{100}{\micro\second}~: un temporisateur va dans la case $\lfloor t/100\,\mu\text{s}\rfloor \bmod 1\,024$,
et un temporisateur plus éloigné que l'étendue de la roue, soit environ \qty{0.1}{\second}, reste simplement dans sa case jusqu'à ce que son heure revienne.
Deux détails la rendent déterministe~: les temporisateurs échus dans une case se déclenchent triés par instant puis par identifiant, jamais
dans l'ordre d'insertion~; et un temporisateur programmé pendant qu'un autre se déclenche, s'il est échu dans la même case, se déclenche au même passage.

\begin{omfigure}
\begin{tikzpicture}[font=\scriptsize,>=Stealth]
\foreach \i in {0,...,15} {
  \pgfmathsetmacro{\a}{90 - 22.5*\i}
  \pgfmathtruncatemacro{\busy}{(\i==2)||(\i==3)||(\i==7)||(\i==12)}
  \ifnum\busy=1 \draw[fill=omMeth!25,draw=omMeth] (\a+11.25:1.4) arc (\a+11.25:\a-11.25:1.4) -- (\a-11.25:2.3) arc (\a-11.25:\a+11.25:2.3) -- cycle;
  \else \draw[fill=gray!6,draw=gray!50] (\a+11.25:1.4) arc (\a+11.25:\a-11.25:1.4) -- (\a-11.25:2.3) arc (\a-11.25:\a+11.25:2.3) -- cycle; \fi
  \node at (\a:1.85) {\i};
}
\draw[->,very thick,omThm] (0,0) -- (45:1.35);
\node[align=center] at (0,-0.35) {maintenant~: case 2};
\node[anchor=west,align=left] at (2.9,1.4) {case $= \lfloor t / g \rfloor \bmod 16$ ($g$~: granularité)};
\node[anchor=west,align=left] at (2.9,0.7) {case 3~: acquittement de la cotation 41 à $t_0 + 50\,\mu$s;\\ contrôle de péremption à $t_0 + 20$ ms};
\node[anchor=west,align=left] at (2.9,-0.2) {avancer à $t$~: visiter les cases 2, 3, \dots\ jusqu'à celle de $t$,\\ déclencher les temporisateurs échus avant $t$ dans l'ordre (instant, identifiant)~;};
\node[anchor=west,align=left] at (2.9,-1.1) {un temporisateur à plus d'un tour\\ (le contrôle de péremption) attend dans sa case};
\end{tikzpicture}

\omcaption{Une \omterm{def:ll:the-strategy-engine:wheel}{roue de temporisateurs} à hachage de 16 cases (celle du build en compte 1\,024 de \qty{100}{\micro\second}). Programmer place un temporisateur
dans la case de son instant d'échéance~; avancer visite les cases franchies et ne déclenche que les temporisateurs échus, triés~; un temporisateur situé au-delà
d'un tour reste dans sa case jusqu'à ce que son heure revienne.}\label{fig:ll:the-strategy-engine:wheel}
\end{omfigure}

\omcode{code/firm/stratengine/cpp/firm_stratengine.hpp}{82}{101}{Faire avancer la roue~: visiter les cases franchies, déclencher les temporisateurs
échus, triés par instant puis par identifiant.}\label{lst:ll:the-strategy-engine:wheel}

La roue ne connaît aucune horloge. Elle est avancée à l'instant de chaque événement, et le code du moteur demande à une \omterm{def:ll:the-strategy-engine:wheel}{horloge injectée} l'heure~:
l'horloge simulée rend l'instant de l'événement, l'horloge réelle lit \texttt{CLOCK\_MONOTONIC} (chapitre~5) et la projette sur la
journée. Rejouer une journée avec l'horloge simulée donne, événement par événement, ce que le moteur en direct aurait fait avec les mêmes entrées
aux mêmes instants~; l'horloge réelle ne sert qu'en direct, et seulement là.

\omcode{code/firm/stratengine/cpp/firm_stratengine.hpp}{46}{58}{Les deux horloges que le moteur peut recevoir.}\label{lst:ll:the-strategy-engine:clock}

\section{Configuration et paramètres à l'exécution}

\begin{definition}[Instantané de paramètres]\label{def:ll:the-strategy-engine:snapshot}
Un \emph{instantané de paramètres} est un ensemble complet et versionné des paramètres d'une stratégie, qui prend effet à un instant déclaré, entre
deux événements et jamais pendant un événement~; chaque action que la stratégie accomplit enregistre la version des paramètres sous lesquels elle a été accomplie.
\end{definition}

Les traders modifient les paramètres au cours de la journée~: une fourchette plus large avant une annonce, une taille plus petite sur un marché mince. Modifier un champ d'une
structure vivante depuis un autre fil pendant que le moteur la lit est une \omterm{def:ll:rust-for-low-latency:race}{course aux données} (data race, chapitre~11), et même une mise à jour atomique correcte d'un
seul champ laisse un instant où la moitié des paramètres sont neufs. Le build applique des instantanés entiers, dans la boucle, entre les événements~:
le chemin de contrôle publie un instantané assorti d'un instant d'effet sur l'anneau de contrôle du moteur, et la boucle l'installe avant le
premier événement à cet instant ou après. La version estampillée sur chaque action indique quels paramètres l'ont produite, et en rejeu les
instantanés sont des entrées comme les autres événements~: le test du build en applique un à mi-journée et vérifie que chaque action
antérieure est inchangée et que chaque action ultérieure porte la nouvelle version.

\section{Discipline du chemin critique}

Le chemin critique (hot path) du chapitre~1 traverse ce fil, et il est discipliné par des règles que les chapitres précédents ont mesurées. Aucune
allocation par événement (chapitre~6)~: le test du build compte les allocations de tas pendant que le moteur rejoue sa fixture et n'en trouve aucune,
parce que le journal des actions, les cases de la roue et le carnet sont dimensionnés au démarrage. Aucun appel système~: l'horloge simulée est un champ,
et même l'horloge réelle se lit par le vDSO (chapitre~5). Aucun verrou ni aucune attente~: les entrées et les sorties sont des anneaux, et un
anneau de sortie plein est une erreur à signaler, non une raison d'attendre. Aucune mise en forme de texte~: le journal reçoit des enregistrements binaires (chapitre~23). Aucun
travail non borné dans un gestionnaire~: le travail de la roue est borné par les cases franchies, celui du carnet par sa structure.

Le résultat sur le portable~: le moteur prend environ \qty{53}{\nano\second} par événement en médiane sur le flux à grand tick du
chapitre~18, mise à jour du carnet comprise, et environ \qty{57}{\nano\second} sur celui à petit tick où il recote à presque chaque
événement~; le p99 est d'environ \qty{130}{\nano\second} et \qty{200}{\nano\second}. Une décision en dizaines de nanosecondes laisse le
\omterm{def:ll:where-latency-comes-from:budget}{budget de latence} au réseau et à la passerelle.

\section{Tutoriel~: rejouer une journée, puis la casser}\label{tut:ll:the-strategy-engine:tut}

\begin{tutorial}
\textbf{Objectif.} Rejouer une entrée enregistrée dans le moteur jusqu'à obtenir la même sortie à chaque fois, puis activer chaque source de
non-déterminisme et compter les issues. \textbf{État final~:} la \cref{fig:ll:the-strategy-engine:reruns}, le temps du moteur par
événement, et des tests au vert en C++ et en Rust.
\begin{tutsteps}
\item \textbf{Rejeu.} Le test C++ fait tourner le moteur deux fois sur la fixture à petit tick du \omterm{def:ll:the-order-book-builder:builder}{constructeur de carnet} et compare les
empreintes de sortie, les actions et le profit, puis les confronte à \texttt{data/expected.txt}~; le moteur Rust, écrit selon les
mêmes règles, reproduit la même empreinte.
\item \textbf{Modifier les paramètres à mi-parcours} avec un instantané, et vérifier que seules les actions qui suivent changent.
\item \textbf{Compter les zéros~:} le test compte les allocations de tas pendant le rejeu.
\item \textbf{Casser le système} avec \texttt{python bench\_engine.py}~: cent nouveaux processus par configuration de \texttt{ll\_nondet},
chacun affichant son empreinte de sortie, puis les empreintes distinctes sont comptées.
\end{tutsteps}
\textbf{Pour aller plus loin.} Remplacer l'ensemble de hachage à graine du petit moteur par un ensemble à graine fixe et recompter~; faire passer les
événements des deux fils par une fusion par horodatage au lieu d'une file partagée.
\end{tutorial}

\section{Construction~: le moteur de stratégie}\label{bld:ll:the-strategy-engine:stratengine}

\begin{build}
\bfield{Rôle} Le fil sur lequel tournent les stratégies de la firme~: il consomme la vue du \omterm{def:ll:the-order-book-builder:builder}{constructeur de carnet} (chapitre~19), envoie des ordres
par la passerelle (chapitre~21) en passant la barrière de risque (chapitre~22), et journalise tout ce qu'il fait (chapitre~23).
\bfield{Interface} C++20 \texttt{firm::strat}~: \texttt{Engine<Clock>(tick, Params, clock)} avec \texttt{run(events, snapshots)},
\texttt{actions}, \texttt{hash}, \texttt{position}, \texttt{cash}, \texttt{pnl()}~; \texttt{TimerWheel(granularity, slots)} avec
\texttt{schedule} et \texttt{advance}~; \texttt{SimClock}, \texttt{RealClock}~; \texttt{Params}, \texttt{Snapshot},
\texttt{Action}. Rust \texttt{firm\_stratengine}~: \texttt{Engine}, \texttt{TimerWheel}, \texttt{Params}, \texttt{Action},
\texttt{read\_events}.
\bfield{Règles} Un seul fil possède l'état~; chaque événement est traité jusqu'à complétion~; les temporisateurs échus avant un événement se déclenchent avant lui, dans l'ordre (instant,
identifiant)~; les instantanés de paramètres s'appliquent entre les événements~; le temps vient de l'\omterm{def:ll:the-strategy-engine:wheel}{horloge injectée}~; la stratégie d'exemple calcule
en entiers~; aucune allocation par événement.
\bfield{Tests d'acceptation} \texttt{code/firm/stratengine/}~: deux rejeux aux actions, à l'empreinte et au profit identiques, égaux
à l'empreinte validée~; le moteur Rust égal à celle-ci~; un instantané ne modifiant que ce qui le suit~; zéro allocation pendant le rejeu~; la
roue se déclenchant dans l'ordre à travers les cases et les tours.
\bfield{Pour aller plus loin} Une roue hiérarchique pour les longs temporisateurs~; des rapports d'ordres issus du simulateur du livre~10 à la place du
bouchon de la plateforme~; plusieurs stratégies sur un même moteur avec un ordre de programmation fixe.
\end{build}

\begin{omsources}
\item G.~Varghese et T.~Lauck, ``Hashed and hierarchical timing wheels: data structures for the efficient implementation of a
timer facility'', \emph{Proceedings of the 11th ACM Symposium on Operating Systems Principles}, 1987.
\item M.~Fowler, \emph{The LMAX Architecture}, 2011.
\item The Rust standard library, \texttt{std::collections::HashMap}.
\end{omsources}

\section{Exercices}

\begin{exercise}[$\star$]\label{exo:ll:the-strategy-engine:1}
Une roue comporte 1\,024 créneaux de \qty{100}{\micro\second}. Dans quel créneau tombe un temporisateur échéant à 34\,200.012345 secondes après minuit,
et à combien de révolutions se trouve un temporisateur échéant dans 5 secondes~?
\end{exercise}

\begin{exercise}[$\star$]\label{exo:ll:the-strategy-engine:2}
Trois temporisateurs échéent dans le même créneau à 120, 121 et 120 microsecondes, avec les identifiants 7, 3 et 9. Dans quel ordre se déclenchent-ils~?
\end{exercise}

\begin{exercise}[$\star$]\label{exo:ll:the-strategy-engine:3}
Pourquoi l'instant d'un événement est-il le bon «~maintenant~» pour une stratégie en rejeu, et que doit-il être vrai du système
en production pour que le rejeu lui corresponde~?
\end{exercise}

\begin{exercise}[$\star\star$]\label{exo:ll:the-strategy-engine:4}
Si un résultat dépend de l'un des $k$ ordres de parcours également probables qu'un ensemble à hachage peut produire,
combien de résultats distincts une centaine de relances devrait-elle montrer en moyenne~? Que suggère le 93 mesuré quant à $k$~?
\end{exercise}

\begin{exercise}[$\star\star$]\label{exo:ll:the-strategy-engine:5}
Un trader fait passer la fourchette de deux ticks à quatre à 10:30:00.000. Un événement arrive à 10:29:59.999 et le suivant à
10:30:00.004. Sous quels paramètres chacun est-il traité, et que montre le journal~?
\end{exercise}

\begin{exercise}[$\star\star$]\label{exo:ll:the-strategy-engine:6}
Le moteur consomme \qty{53}{\nano\second} par événement à la médiane. À l'ouverture, les événements arrivent à 2 millions par seconde pendant 50
millisecondes. Suit-il le rythme, et quel est son taux d'utilisation~?
\end{exercise}

\begin{exercise}[$\star\star\star$]\label{exo:ll:the-strategy-engine:7}
\emph{Programmation.} Ajouter au moteur un temporisateur de rafraîchissement~: chaque cotation vieille de plus de \qty{5}{\second} est renvoyée au même prix. Vérifier
que deux rejeux concordent toujours et que le moteur Rust, modifié de la même façon, produit toujours le hachage du moteur C++.
\end{exercise}

\begin{exercise}[$\star\star\star$]\label{exo:ll:the-strategy-engine:8}
\emph{Trouver la faille.} «~Notre moteur est déterministe~: il est mono-thread et n'utilise aucun aléa. Il estampille chaque ordre avec
\texttt{std::chrono::system\_clock::now()} pour la piste d'audit.~»
\end{exercise}

\section{Problème~: Deux exécutions, deux profits}

\begin{problem}[{Problème du week-end --- compter les façons dont un programme peut être en désaccord avec lui-même}]\label{pb:ll:the-strategy-engine:1}
Une équipe de recherche rejoue la même journée enregistrée dans son moteur de cotation à deux reprises et obtient deux profits. Utiliser les mesures de la
\cref{fig:ll:the-strategy-engine:reruns} (\texttt{measured\_nondet.csv}).

\medskip\noindent\textbf{Partie I --- Les sources.}
\begin{enumerate}
\item Lister les trois sources mesurées et le nombre de résultats distincts que chacune a donnés sur une centaine d'exécutions.
\item Pourquoi l'horloge murale en donne-t-elle une centaine~?
\item Pourquoi l'ensemble à hachage initialisé avec une graine en donne-t-il moins d'une centaine~?
\item Pourquoi le flux à deux threads en donne-t-il encore moins, et qu'est-ce qui pourrait en donner davantage~?
\end{enumerate}

\medskip\noindent\textbf{Partie II --- Les repérer.}
\begin{enumerate}[resume]
\item Comment on trouverait la première action à laquelle deux exécutions divergent~?
\item Que laisse chacune des sources comme indice à ce point~?
\item Pourquoi chaque relance doit-elle être un nouveau processus pour faire apparaître l'effet du hachage initialisé avec une graine~?
\item Laquelle des trois sources ne se manifesterait jamais dans une seule exécution sous débogueur~?
\end{enumerate}

\medskip\noindent\textbf{Partie III --- Les correctifs.}
\begin{enumerate}[resume]
\item Corriger l'horloge murale.
\item Corriger l'ordre de parcours du hachage sans renoncer aux tables de hachage.
\item Corriger le flux à deux threads.
\item Après les correctifs, combien de résultats distincts une centaine de relances devraient-elles donner, et combien le moteur de la société en a-t-il donné~?
\end{enumerate}

\medskip\noindent\textbf{Partie IV --- Le verdict.}
\begin{enumerate}[resume]
\item Énoncer le \emph{résultat nommé}~: le nombre de résultats distincts sur une centaine de relances avec chaque source activée, puis après
les correctifs.
\item Pourquoi «~le profit diffère légèrement~» n'est-il pas une raison d'accepter le non-déterminisme dans un backtest~?
\item Qu'est-ce qui peut encore rompre la reproductibilité bit à bit entre deux machines, même avec les trois correctifs~?
\item Pourquoi l'arithmétique entière de la stratégie fait-elle concorder les moteurs C++ et Rust~?
\item Que doit exécuter le pipeline d'intégration continue pour préserver le déterminisme~?
\item Qu'apporte le déterminisme à l'équipe de support de production~?
\item Quand un composant non déterministe est-il acceptable dans un système de trading~?
\item En une phrase~: qu'est-ce qui rend un moteur rejouable~?
\end{enumerate}
\end{problem}

\section{Questions d'entretien}

\begin{interviewq}[$\star$ \iqroles{développeur}]\label{iq:ll:the-strategy-engine:1}
Pourquoi exécuter une stratégie sur un seul thread~? Quel en est le coût~?
\end{interviewq}

\begin{interviewq}[$\star\star$ \iqroles{développeur}]\label{iq:ll:the-strategy-engine:2}
Comment implémenteriez-vous des temporisateurs pour des milliers d'ordres, chacun avec ses propres délais d'expiration~?
\end{interviewq}

\begin{interviewq}[$\star\star$ \iqroles{développeur}]\label{iq:ll:the-strategy-engine:3}
Votre backtest et votre exécution en production divergent le même jour. Où regardez-vous~?
\end{interviewq}

\begin{interviewq}[$\star\star$ \iqroles{développeur}]\label{iq:ll:the-strategy-engine:4}
Comment changez-vous les paramètres d'une stratégie pendant qu'elle tourne, en toute sécurité~?
\end{interviewq}

\begin{interviewq}[$\star\star$ \iqroles{développeur}]\label{iq:ll:the-strategy-engine:5}
Que ne doit jamais se produire sur le chemin critique (hot path) d'un moteur de stratégie, et comment le vérifie-t-on~?
\end{interviewq}

\begin{interviewq}[$\star\star\star$ \iqroles{développeur, chercheur}]\label{iq:ll:the-strategy-engine:6}
Concevez un moteur de stratégie à la fois en production et rejouable bit à bit~: quelles sont les règles, et comment les fait-on appliquer~?
\end{interviewq}
